\documentclass[12pt,a4paper]{article}

\usepackage{styles/couleurs}
\usepackage{styles/mise_en_page}
\usepackage{styles/boites}
\usepackage{styles/maths}
\usepackage{styles/figures}

\usepackage{enumitem}
\usepackage{array}
\usepackage{tabularx}
\usepackage{amssymb}
\usepackage{tikz}

\setlist[enumerate]{leftmargin=*,itemsep=0.35em,topsep=0.35em}

\newcommand{\lignes}[1]{%
  \begin{tcolorbox}[
    colback=white,
    colframe=NeutreC!40,
    boxrule=0.5pt,
    arc=1.5mm,
    left=3mm,right=3mm,top=2mm,bottom=1mm
  ]
    \foreach \i in {1,...,#1}{\noindent\dotfill\par\vspace{0.45em}}
  \end{tcolorbox}
}

\newcommand{\blocdonnees}[1]{%
  \begin{tcolorbox}[
    mathsbox,
    colback=NeutreBg,
    colframe=NeutreC!45,
    boxrule=0.6pt,
    arc=2mm
  ]
  #1
  \end{tcolorbox}
}

\begin{document}

\begin{center}
  {\sffamily\LARGE\bfseries\color{TitreC}Séquence 1 -- Statistiques}\\[0.2em]
  {\sffamily\large\bfseries\color{SavoirC}Activité 1 -- Bowling}
\end{center}

\blocdonnees{
Treize élèves jouent une partie de bowling.

\medskip
\textbf{Équipe des Verts -- 7 joueurs}
\[
105 \qquad 120 \qquad 104 \qquad 121 \qquad 99 \qquad 127 \qquad 108
\]

\textbf{Équipe des Bleus -- 6 joueurs}
\[
93 \qquad 181 \qquad 89 \qquad 98 \qquad 117 \qquad 94
\]
}

\begin{enumerate}
  \item Calculer puis interpréter la moyenne des deux équipes.

  \medskip
  \noindent\textbf{Moyenne des Verts =}

  \vspace{1.1em}

  \noindent\textbf{Moyenne des Bleus =}

  \lignes{4}

  \item Les élèves de l'équipe verte étudient leurs résultats.
  Après avoir calculé la moyenne, ils utilisent le nombre \(108\)
  pour parler de leur performance.

  \medskip
  \textbf{Comment ont-ils choisi ce nombre ?}

  \lignes{7}

  \item Calculer la médiane de l'équipe bleue.

  \lignes{5}
\end{enumerate}

% NOTES PROFESSEUR -- NON IMPRIMÉES
% Activité volontairement concise.
% Relances graduées possibles :
% 1. « La présentation actuelle des scores vous aide-t-elle ? »
% 2. « Peut-on écrire les scores autrement pour mieux voir ce qui se passe ? »
% 3. « Essaie de les ranger. »

\end{document}
